% ------------------------------------------------------------------
%  Capitolo 4 — Da Simonide a Ebbinghaus
%  Parte II
%  Ricerca di riferimento: ricerca 01 §2; 03 (intero)
%  Voci bibliografiche principali: ebbinghaus1885, james1890, spitzer1939, carruthers1990, yates1966, rossi1960, bolzoni1995
%  Epigrafe: ✔ verificata sul testo (vedi ricerca 03 e registro, ricerca 00)
%  Scheda del capitolo: ricerca/06_indice-ragionato.md, capitolo 4
% ------------------------------------------------------------------
\chapter{Da Simonide a Ebbinghaus}
\label{cap:storia}

\epigrafe[è soprattutto l'ordine a dare luce alla memoria]{ordinem esse maxime, qui memoriae lumen adferret}{Cicerone, De oratore II, 353}

\iniziale{Q}{uesto capitolo} \segnaposto{apertura: da scrivere — vedi la scheda nell'indice ragionato}.

\section{Il banchetto di Scopas}

\segnaposto{da scrivere}

\section{Aristotele e l'arte di ricordare}

\segnaposto{da scrivere}

\section{Retori, monaci e teologi}

\segnaposto{da scrivere}

\section{La nascita di una scienza}

\segnaposto{da scrivere}

\begin{daricordare}
\segnaposto{il principio del capitolo in due righe}
\end{daricordare}

\begin{domande}
  \item \segnaposto{domanda di recupero su questo capitolo}
  \item \segnaposto{domanda di applicazione a una materia che studi}
  \item \segnaposto{domanda di ripresa su un capitolo precedente}
\end{domande}

\begin{sintesi}
  \item \segnaposto{punto di sintesi}
\end{sintesi}

\finecapitolo
